% Part: first-order-logic
% Chapter: introduction
% Section: sentences

\documentclass[../../../include/open-logic-section]{subfiles}

\begin{document}

\olfileid{fol}{int}{snt}

\olsection{\usetoken{P}{sentence}}

We beschikken nu over een schets van een definitie van vervulling
(``waar in'') voor !!{structure}s en !!{formula}s. Daarvoor is echter
een extra gegeven nodig: !!a{variable} bedeling. Wij wilden juist een
definitie van !!{sentence}s geven. Hoe kunnen we de bedelingen
elimineren, en wat zijn !!{sentence}s?

Men herinnert zich wellicht uit een eerste inleiding in de formele
logica de bespreking van het verband tussen !!{variable}s en
kwantoren. Op een kwantor volgt steeds !!a{variable}. In het deel van
de !!{sentence} waarop die kwantor betrekking heeft, zijn
``bereik'', wordt die !!{variable} door de kwantor ``gebonden''
geacht. Bij !!{formula}s is niet vereist dat iedere !!{variable} een
bijbehorende kwantor heeft; !!{variable}s zonder bijbehorende kwantor
zijn ``vrij'' of ``ongebonden''. We nemen als !!{sentence}s alle
!!{formula}s die geen vrije !!{variable}s hebben.

Het intuïtieve idee wanneer een voorkomen van !!a{variable} in
!!a{formula}~$!A$ gebonden is, door welke kwantor het wordt gebonden en
wanneer het vrij is, is wederom niet moeilijk te doorzien. Mogelijk
heeft men een toetsmethode geleerd, bijvoorbeeld door haakjes te
tellen. We moeten echter een precieze definitie geven. Omdat we
!!{formula}s inductief hebben gedefinieerd, kunnen we ook de vrije en
gebonden voorkomens van !!a{variable}~$x$ in !!a{formula}~$!A$
inductief definiëren. Voor onze vereenvoudigde taal zou die definitie
er bijvoorbeeld als volgt uitzien:
\begin{enumerate}
  \item Als $!A$ atomair is, zijn alle voorkomens van $x$ daarin vrij
  (dus is het voorkomen van $x$ in~$\Atom{\Obj P}{x}$ vrij).
  \item Als $!A$ de vorm $\lnot !B$ heeft, is een voorkomen van~$x$
  in~$\lnot !B$ vrij dan en slechts dan als het overeenkomstige
  voorkomen van~$x$ in~$!B$ vrij is (de vrije voorkomens van variabelen
  in~$!B$ zijn dus precies de overeenkomstige voorkomens in~$\lnot !B$).
  \item Als $!A$ de vorm $(!B \land !C)$ heeft, is een voorkomen van~$x$
  in~$(!B \land !C)$ vrij dan en slechts dan als het overeenkomstige
  voorkomen van~$x$ in~$!B$ of in~$!C$ vrij is.
  \item Als $!A$ de vorm $\lexists[x][!B]$ heeft, is geen enkel voorkomen
  van $x$ in~$!A$ vrij. Heeft het de vorm $\lexists[y][!B]$, waarbij
  $y$ een andere !!{variable} is dan~$x$, dan is een voorkomen van~$x$
  in $\lexists[y][!B]$ vrij dan en slechts dan als het overeenkomstige
  voorkomen van~$x$ in~$!B$ vrij is.
\end{enumerate}

Zodra we een precieze definitie van vrije en gebonden voorkomens van
variabelen hebben, kunnen we eenvoudig stellen: !!a{sentence} is elke
!!{formula} zonder vrije voorkomens van !!{variable}s.

\end{document}
